%==============================================================================
\section{How to read this guide}
%==============================================================================

\paragraph{What FastHydrology computes.} An ice-flow model needs the effective pressure $N = p_o - p_w$ at the
bed (the ice overburden minus the water pressure), because $N$ sets how fast the ice slides. FastHydrology
computes $N$, and a water-layer thickness $W$, on a regular 2-D map-view grid, from the ice geometry and the
basal water supply. Its own models are steady: the water system is assumed to adjust much faster than the ice, so
each call recomputes $N$ from the present ice state, with no memory of the previous call. This makes them cheap
enough to call once per ice-model time step, or even inside the ice model's velocity iteration.

\paragraph{The three models.}
\begin{itemize}[leftmargin=1.5em]
\item \code{KazmierczakHydroModel} (``K24'', Kazmierczak et al.\ 2024, \emph{The Cryosphere} 18, 5887). The basal
  melt and any water from above are routed down the hydraulic potential to give a distributed flux $q$
  (Le Brocq et al.\ 2006, 2009). From $q$ and the local potential gradient, a conduit balance gives a far-field
  effective pressure $\Ninf$, and an error-function transition brings $N$ to zero at the grounding line. The bed
  can be hard, soft or mixed, and the drainage switches between efficient and inefficient by itself. Nearly all
  of the code is this model.
\item \code{HABHydroModel} (``HAB'', Kazmierczak et al.\ 2022, Sec.~2.1.1). The water pressure is a fixed fraction
  of the ocean pressure at the bed: a closed formula, a few lines of code.
\item \code{ShaktiHydroModel}. A wrapper around a \code{Shakti.Simulation} from the separate Shakti.jl package
  (Sommers et al.\ 2018), a time-evolving model. FastHydrology only forwards \code{run!} and \code{step!} to it.
\end{itemize}

\paragraph{How the guide is organised.}
\begin{itemize}[leftmargin=1.5em]
\item Section~\ref{sec:physics} states the equations of K24 and HAB with the symbols the code uses.
\item Section~\ref{sec:discretisation} shows how they become array operations: the grid, the gradients, pit
      filling, the routing graph, the conversion of routed flux to $q$, and the stress-gradient smoothing.
\item Section~\ref{sec:architecture} explains how the package is built: the grid interface, options as types,
      the split of the K24 model into parameters and workspace, the extensions, and the load order.
\item Section~\ref{sec:lifecycle} follows one simulation from \code{ArrayHydroGrid} to \code{run!}, and one
      steady-state solve through every function it calls. Read this section to understand the code's flow.
\item Section~\ref{sec:reference} documents each source file, every function in it, what it computes and why it
      is written the way it is. Read this section when working on a file.
\item Sections~\ref{sec:routing}--\ref{sec:fixedpoints} go deeper into the two most intricate parts: the routing
      algorithms, and the fixed-point iterations of the dissipation melt and of $(q, N)$.
\item Sections~\ref{sec:options}--\ref{sec:pitfalls} collect every option and default, the coupling interface,
      the tests, the known pitfalls and the open issues. Appendix~\ref{app:fields} lists all fields of the K24
      model; Appendix~\ref{app:index} indexes every exported name.
\end{itemize}

\paragraph{The companion script.} \code{test/code\_tour/code\_tour.jl} is a narrated, runnable walk through the
same material. PART $n$ tours one source file: it prints what the file is for, calls every definition in it on a
small fixture, and checks the result against something independent (a hand calculation, an invariant such as
conservation of routed water, or a second code path that must agree). A coverage summary at the end parses the
source files at run time and reports any definition the tour did not exercise. Each reference subsection here
starts with ``Tour PART $n$''. At \code{c88bd23} it runs 271 checks, all passing, and exercises all 191 definitions of
\file{src/} and \file{ext/} in about two minutes. The same material is in \code{code\_tour.ipynb}, generated from the
script by \code{make\_notebook.py} (edit the script, not the notebook); the job below also runs the notebook's code cells
in order (\code{notebook\_extracted\_check.jl}). On albedo:
\begin{verbatim}
sbatch ~/build_jobs/fh_tour.sbatch     # -> ~/build_jobs/fh_tour_<jobid>.out
\end{verbatim}
It has its own environment (\code{test/code\_tour/Project.toml}: FastHydrology and Shakti developed from the
local repositories, plus CairoMakie, NCDatasets, MAT, Preferences), because CairoMakie and Shakti are only weak
dependencies of FastHydrology. Its \code{LocalPreferences.toml} sets Shakti to the CPU without CUDA. The last
PART runs both models on the real Thwaites 2\,km and Antarctica 16\,km inputs in
\code{test/Kazmierczak et al 2024/input/}. The older three-part notebook tour is kept in
\code{test/code\_tour/previous/}.

\paragraph{Conventions.} Code names are in \code{typewriter}. Fields are \code{Nx}\,$\times$\,\code{Ny} arrays
indexed \code{[i, j]}, $i$ along $x$, $j$ along $y$. ``Grounded'' means \code{mask[i, j] == 1} exactly; the hydrology
is only defined there. A file name like \file{common/grid.jl} is relative to the repository root. Units are SI
throughout the code (m, s, Pa, kg, W): velocities and rate factors given per year must be converted with
\code{perYear2perSecond} before they reach a model, and the data loaders do this. ``K24 Eq.~$n$'' refers to
Kazmierczak et al.\ (2024); ``Le Brocq Eq.~$n$'' to Le Brocq, Payne and Siegert (2006, \emph{Computers \&
Geosciences} 32, 1780).

\paragraph{Where to find more.} The Documenter site (\code{docs/}) has the public API and four Literate examples
(Kazmierczak 2024 on Thwaites and Yelmo restarts, HAB, Shakti coupling, a synthetic domain). The coupled
Yelmo--FastHydrology drivers are in Kryonomos.jl (\code{examples/fasthydrology/}). The companion guide for
Shakti.jl is \code{Shakti.jl/test/code\_guide/shakti\_code\_guide.tex}.
